\section{Objectif et programme réalisé}
L'objectif est de comparer trois détecteurs sur un lot d'images et de mesurer leur accord avec des contours tracés à la main. Le programme, écrit en \textbf{C++17 avec OpenCV}, propose une interface à curseurs pour choisir l'image, le détecteur et ses paramètres, ainsi qu'un mode par lot qui exporte les cartes de contours et les mesures au format CSV.

La base de Groningen contient \textbf{40 images de $512\times512$ pixels et 40 références}~\cite{groningen}. Dix images servent au réglage, puis les paramètres sont figés pour les trente autres. Les mesures combinent une comparaison pixel à pixel et une comparaison tolérant un décalage d'un pixel.
